% ============================================================
%  Manual do Usuário — Remote Firmware Flasher v2.0
%  Compilar com: tectonic manual.tex
% ============================================================
\documentclass[12pt,a4paper]{article}

% --- Codificação e língua ----------------------------------------
\usepackage[utf8]{inputenc}
\usepackage[T1]{fontenc}
\usepackage[brazil]{babel}

% --- Tipografia --------------------------------------------------
\usepackage{lmodern}
\usepackage{microtype}

% --- Layout ------------------------------------------------------
\usepackage[
  top=2.5cm, bottom=2.5cm,
  left=2.8cm, right=2.8cm,
  headheight=15pt
]{geometry}
\usepackage{fancyhdr}
\usepackage{titlesec}
\usepackage{parskip}
\setlength{\parindent}{0pt}
\setlength{\parskip}{6pt}
\usepackage{enumitem}

% --- Cor e links -------------------------------------------------
\usepackage{xcolor}
\definecolor{azulescuro}{RGB}{15,52,120}
\definecolor{azulmedio}{RGB}{40,100,180}
\definecolor{azulclaro}{RGB}{220,235,250}
\definecolor{cinzatexto}{RGB}{80,80,80}
\definecolor{amarelodest}{RGB}{255,243,205}
\definecolor{amarelob}{RGB}{180,130,0}
\definecolor{verdedest}{RGB}{220,245,220}
\definecolor{verdeb}{RGB}{40,120,40}

\usepackage[
  colorlinks=true,
  linkcolor=azulmedio,
  urlcolor=azulmedio,
  citecolor=azulmedio
]{hyperref}

% --- Tabelas -----------------------------------------------------
\usepackage{booktabs}
\usepackage{array}
\usepackage{colortbl}
\usepackage{tabularx}

% --- Caixas e destaques ------------------------------------------
\usepackage{tcolorbox}
\tcbuselibrary{skins,breakable}

\newtcolorbox{notabox}[1][Nota]{
  colback=amarelodest, colframe=amarelob,
  fonttitle=\bfseries\small, title=#1,
  left=4pt, right=4pt, top=3pt, bottom=3pt,
  breakable
}

\newtcolorbox{dicabox}[1][Dica]{
  colback=verdedest, colframe=verdeb,
  fonttitle=\bfseries\small, title=#1,
  left=4pt, right=4pt, top=3pt, bottom=3pt,
  breakable
}

\newtcolorbox{codigobox}{
  colback=gray!10, colframe=gray!40,
  fontupper=\ttfamily\small,
  left=8pt, right=4pt, top=4pt, bottom=4pt,
  breakable
}

% --- Código ------------------------------------------------------
\usepackage{listings}
\lstset{
  basicstyle=\ttfamily\footnotesize,
  keywordstyle=\color{azulmedio}\bfseries,
  commentstyle=\color{gray},
  stringstyle=\color{verdeb},
  backgroundcolor=\color{gray!8},
  frame=single,
  framerule=0.4pt,
  rulecolor=\color{gray!40},
  breaklines=true,
  showstringspaces=false,
  tabsize=2
}

\lstdefinelanguage{OIL}{
  morekeywords={OIL_VERSION, CPU, OS, STATUS, BUILD, TRUE, FALSE, STANDARD,
    TRAMPOLINE_BASE_PATH, APP_NAME, APP_SRC, CPPCOMPILER, COMPILER, LINKER,
    ASSEMBLER, COPIER, SYSTEM, PYTHON, LIBRARY, SYSTEM_CALL, APPMODE,
    ALARM, COUNTER, SystemCounter, ACTION, ACTIVATETASK, TASK, AUTOSTART,
    ALARMTIME, CYCLETIME, PRIORITY, ACTIVATION, SCHEDULE, FULL, STACKSIZE},
  morecomment=[l]{//},
  morestring=[b]",
  sensitive=true
}

% --- Cabeçalho e rodapé ------------------------------------------
\pagestyle{fancy}
\fancyhf{}
\fancyhead[L]{\small\color{cinzatexto}Remote Firmware Flasher v2.0}
\fancyhead[R]{\small\color{cinzatexto}Manual do Usuário}
\fancyfoot[C]{\small\color{cinzatexto}\thepage}
\renewcommand{\headrulewidth}{0.4pt}
\renewcommand{\headrule}{\hbox to\headwidth{\color{azulmedio}\leaders\hrule height \headrulewidth\hfill}}

% --- Títulos de seção --------------------------------------------
\titleformat{\section}
  {\Large\bfseries\color{azulescuro}}
  {\thesection.}{0.6em}{}
  [\color{azulmedio}\rule{\linewidth}{1.2pt}]

\titleformat{\subsection}
  {\large\bfseries\color{azulmedio}}
  {\thesubsection}{0.6em}{}

\titleformat{\subsubsection}
  {\normalsize\bfseries\color{azulmedio}}
  {\thesubsubsection}{0.6em}{}

% ================================================================
\begin{document}

% --- Capa --------------------------------------------------------
\begin{titlepage}
  \pagecolor{azulescuro}
  \color{white}
  \vspace*{2cm}

  \begin{center}
    {\fontsize{11}{13}\selectfont\scshape Residência Tecnológica --- UFPE/CIn}\\[0.8cm]
    \rule{\linewidth}{1.5pt}\\[0.6cm]
    {\fontsize{28}{34}\selectfont\bfseries Remote Firmware\\Flasher}\\[0.4cm]
    \rule{\linewidth}{1.5pt}\\[0.8cm]
    {\fontsize{15}{18}\selectfont Manual do Usuário}\\[0.3cm]
    {\fontsize{12}{15}\selectfont\color{azulclaro}%
      Sistemas Operacionais de Tempo Real}\\[2.5cm]
  \end{center}

  \begin{tcolorbox}[
    colback=white, colframe=white,
    coltext=azulescuro,
    width=0.75\textwidth,
    center, arc=3pt,
    left=10pt, right=10pt, top=6pt, bottom=6pt
  ]
    \small
    \begin{tabular}{@{}ll@{}}
      \textbf{Versão}          & 2.0 \\
      \textbf{Data}            & Março de 2026 \\
      \textbf{Repositório}     & \texttt{renatosfagundes/remote-flasher} \\
      \textbf{Plataforma}      & Windows 10/11 \\
    \end{tabular}
  \end{tcolorbox}

  \vfill
  {\small\itshape Gerado por \texttt{manual.tex} --- Compilado com Tectonic}
\end{titlepage}
\pagecolor{white}
\color{black}

% --- Sumário -----------------------------------------------------
\tableofcontents
\newpage

% === 1. Introdução ===============================================
\section{Introdução}

O \textbf{Remote Firmware Flasher} é uma aplicação desktop (PySide6) que permite programar placas Arduino remotamente, através de conexão VPN e SSH. Foi desenvolvido para a Residência Tecnológica do CIn/UFPE, permitindo que alunos gravem firmware nas placas do laboratório sem precisar estar fisicamente presente.

\subsection{Funcionalidades}

\begin{itemize}[nosep]
  \item Conexão VPN (SSTP) diretamente pela aplicação
  \item Upload e gravação de firmware (\texttt{.hex}) via SCP + \texttt{avrdude}
  \item Reset das placas via scripts PowerShell remotos
  \item Até 4 terminais serial simultâneos com envio de comandos
  \item Terminal SSH para executar comandos no PC remoto
  \item Upload de arquivos e pastas via SFTP com barra de progresso
  \item Câmera ao vivo para visualizar as placas
  \item Configuração automática do ambiente de desenvolvimento (Setup)
  \item Guias de integração com Notepad++ e VSCode
\end{itemize}

% === 2. Instalação ===============================================
\section{Instalação}

\subsection{Executável (recomendado)}

Basta baixar o arquivo \texttt{RemoteFlasher.exe} e executar. Não é necessário instalar Python ou qualquer dependência.

\subsection{Código-fonte}

Para executar a partir do código-fonte:

\begin{codigobox}
git clone https://github.com/renatosfagundes/remote-flasher.git\\
cd remote-flasher\\
pip install -r requirements.txt\\
cp secrets.example.py secrets.py   \# editar com as credenciais do lab\\
python src/main.py
\end{codigobox}

\textbf{Dependências:} Python 3.9+, PySide6, paramiko, requests.

\subsection{Configurando as credenciais (\texttt{secrets.py})}

O arquivo \texttt{secrets.py} contém os endereços IP, usuários e senhas SSH dos PCs do laboratório. Esse arquivo \textbf{não é incluído no repositório} por segurança. Para obtê-lo, entre em contato com \textbf{Davi} ou \textbf{Renato}.

% === 3. Configuração do Ambiente =================================
\section{Configuração do Ambiente de Desenvolvimento}

A aba \textbf{Setup} do aplicativo automatiza toda a configuração do ambiente necessário para compilar e gravar firmware usando o Trampoline RTOS.

\subsection{O que é instalado automaticamente}

\begin{center}
\begin{tabularx}{\textwidth}{lX}
  \toprule
  \textbf{Componente} & \textbf{Descrição} \\
  \midrule
  Arduino CLI       & Ferramenta de linha de comando para compilar e gravar firmware Arduino \\
  AVR Toolchain     & Compilador GCC para microcontroladores AVR (avr-gcc, avr-objcopy, etc.) \\
  Trampoline RTOS   & Sistema operacional de tempo real compatível com OSEK/AUTOSAR \\
  goil              & Compilador OIL (OSEK Implementation Language) do Trampoline \\
  Arduino Core      & Bibliotecas do Arduino para uso com Trampoline \\
  Template Nano     & Template de configuração do Trampoline para Arduino Nano \\
  MCP\_CAN          & Biblioteca para comunicação CAN via módulo MCP2515 \\
  \bottomrule
\end{tabularx}
\end{center}

\subsection{Executando o Setup}

\begin{enumerate}[nosep]
  \item Abra a aba \textbf{Setup}
  \item Os indicadores mostram o status de cada componente (verde = instalado, vermelho = ausente)
  \item Clique em \textbf{``Run Full Setup''} para instalar tudo automaticamente
  \item Acompanhe o progresso no log à esquerda
  \item Após a conclusão, abra um \textbf{novo terminal} para que as variáveis PATH sejam reconhecidas
\end{enumerate}

\begin{notabox}[Importante]
O setup requer conexão com a internet para baixar os componentes. O processo pode levar alguns minutos na primeira execução. Execuções subsequentes são rápidas pois detectam o que já está instalado.
\end{notabox}

\subsection{Verificação manual}

Após o setup, abra um novo terminal e verifique:

\begin{codigobox}
arduino-cli version\\
avr-gcc --version\\
goil --version
\end{codigobox}

\subsection{Estrutura de diretórios}

Todos os componentes são instalados em \texttt{C:\textbackslash ESA}:

\begin{codigobox}
C:\textbackslash ESA\textbackslash\\
\quad arduino\_cli\textbackslash bin\textbackslash\quad\quad\# Arduino CLI\\
\quad avr\_tools\textbackslash\quad\quad\quad\quad\quad\quad\# AVR GCC Toolchain + goil\\
\quad trampoline\textbackslash\quad\quad\quad\quad\quad\# Trampoline RTOS\\
\quad\quad examples\textbackslash avr\textbackslash\quad\quad\# Exemplos (arduinoUno, arduinoNano)\\
\quad\quad goil\textbackslash templates\textbackslash\quad\# Templates de configuração\\
\quad\quad opt\textbackslash devel\textbackslash\quad\quad\quad\# Diretório de desenvolvimento\\
\end{codigobox}

% === 4. Primeira Execução ========================================
\section{Primeira Execução}

Na primeira vez que você abrir o aplicativo, uma janela de configuração será exibida solicitando:

\begin{itemize}[nosep]
  \item \textbf{Seu nome} --- usado para criar sua pasta no PC remoto.
  \item \textbf{Pasta remota} --- preenchida automaticamente como \texttt{c:\textbackslash 2026\textbackslash <seu\_nome>}.
\end{itemize}

Configurações são salvas em \texttt{\%APPDATA\%\textbackslash RemoteFlasher\textbackslash settings.json}.

% === 5. Conectando à VPN ========================================
\section{Conectando à VPN}

A aba \textbf{VPN} permite conectar-se à rede do laboratório:

\begin{enumerate}[nosep]
  \item Insira seu usuário e senha do CIn
  \item (Opcional) Marque ``Remember me'' para salvar as credenciais
  \item Clique em ``Connect VPN''
  \item Aguarde o indicador ficar verde
\end{enumerate}

\begin{notabox}
Se a VPN ainda não estiver configurada no Windows, clique em ``Setup VPN Profile'' primeiro.
\end{notabox}

% === 6. Gravando Firmware ========================================
\section{Gravando Firmware (Flash)}

A aba \textbf{Flash} permite enviar e gravar firmware nas placas.

\begin{enumerate}[nosep]
  \item Selecione o PC de destino (ex: PC 217)
  \item Selecione a placa (ex: Placa 01) e a porta ECU
  \item Clique em ``Browse'' e selecione seu arquivo \texttt{.hex}
  \item Use ``Upload + Reset + Flash'' para executar tudo de uma vez
\end{enumerate}

% === 7. Terminal Serial ==========================================
\section{Terminal Serial}

A aba \textbf{Serial} permite monitorar até 4 conexões seriais simultaneamente e enviar comandos.

\subsection{Layout adaptativo}

\begin{itemize}[nosep]
  \item \textbf{1 conexão:} painel único, área completa
  \item \textbf{2 conexões:} lado a lado (divisor horizontal)
  \item \textbf{3 conexões:} 2 em cima + 1 embaixo
  \item \textbf{4 conexões:} grade 2$\times$2
\end{itemize}

Clique em \textbf{``+ Add Serial''} para adicionar uma nova conexão (máximo 4). Cada painel pode ser fechado pelo botão \textbf{X} (que desconecta automaticamente).

\subsection{Enviando comandos}

Cada painel possui um campo de texto na parte inferior para enviar comandos pela porta serial. Digite o comando e pressione \textbf{Enter} ou clique em \textbf{Send}. O campo só fica ativo quando a conexão está aberta.

\subsection{Autoscroll}

Desmarque a caixa \textbf{``Autoscroll''} para congelar a posição do scroll enquanto os dados continuam chegando --- útil para analisar mensagens CAN sem perder a posição.

\subsection{Exclusão de portas}

Quando um painel está conectado a uma porta COM, essa porta não aparece nos menus dos outros painéis, evitando conflitos.

% === 8. Terminal SSH + Upload SFTP ===============================
\section{Terminal SSH e Upload de Arquivos}

A aba \textbf{SSH Terminal} combina execução de comandos e upload de arquivos.

\subsection{Executar comandos}

Digite um comando e pressione Enter. A saída é exibida no log.

\begin{notabox}
O PC remoto usa Windows. Use comandos como \texttt{dir}, \texttt{cd}, \texttt{type}, e não \texttt{ls}, \texttt{pwd}, \texttt{cat}.
\end{notabox}

\subsection{Upload de arquivos e pastas (SFTP)}

\begin{enumerate}[nosep]
  \item Clique em \textbf{``File...''} para selecionar um arquivo ou \textbf{``Folder...''} para uma pasta inteira
  \item Defina o caminho remoto de destino
  \item Clique em \textbf{``Upload''}
  \item Acompanhe o progresso na barra
\end{enumerate}

A transferência usa SFTP (via paramiko) e cria diretórios remotos automaticamente. Pastas são enviadas recursivamente.

% === 9. Câmera ===================================================
\section{Câmera ao Vivo}

O painel de câmera exibe o feed MJPEG do laboratório. É visível nas abas Flash, Serial e SSH, podendo ser oculto/exibido pelo botão no canto da barra de abas. Com múltiplas conexões serial, a câmera reduz automaticamente para 25\% da largura.

% === 10. Integração com IDEs ====================================
\section{Integração com Notepad++}

A aba \textbf{Setup} contém guias completas de integração. Abaixo está um resumo e um exemplo prático.

\subsection{Plugins necessários}

Instale via \textbf{Plugins $\rightarrow$ Plugins Admin}:
\begin{itemize}[nosep]
  \item \textbf{NppExec} --- execução de comandos e scripts
  \item \textbf{SourceCookifier} --- navegação de código
\end{itemize}

Instale manualmente o \textbf{NppConsole} (baixe de \url{https://sourceforge.net/projects/nppconsole/}).

\subsection{Exemplo: Compilar e gravar um projeto Trampoline}

Crie o seguinte script em \texttt{Plugins $\rightarrow$ NppExec $\rightarrow$ Execute}:

\begin{lstlisting}[caption={Script NppExec: trampoline\_build\_and\_flash.scp}]
// [Trampoline] Build and Flash
NPP_SAVEALL
cd $(CURRENT_DIRECTORY)

// Run goil to generate build files
goil --target=avr/arduino/nano --templates=../../../../goil/templates/ $(NAME_PART).oil

// Compile
cmd /c python make.py

// Flash (change COM port as needed)
INPUTBOX "Serial port:":"COM": "COM25"
set local commport = $(INPUT)
arduino-cli upload -b arduino:avr:nano:cpu=atmega328old -p $(commport) -v -i image.hex -t
\end{lstlisting}

\begin{dicabox}
Para criar um atalho de teclado: \textbf{Plugins $\rightarrow$ NppExec $\rightarrow$ Advanced Options}, associe o script a um item de menu, depois em \textbf{Settings $\rightarrow$ Shortcut Mapper} atribua uma tecla.
\end{dicabox}

\subsection{Máscaras de destaque para erros}

Configure em \textbf{Plugins $\rightarrow$ NppExec $\rightarrow$ Console Output Filters $\rightarrow$ Highlights}:

\begin{center}
\small
\begin{tabular}{clll}
  \toprule
  \textbf{\#} & \textbf{Máscara} & \textbf{Cor} & \textbf{Estilo} \\
  \midrule
  1 & \texttt{\%FILE\%:\%LINE\%:*: error:*} & Vermelho & Negrito \\
  2 & \texttt{\%FILE\%:\%LINE\%:*: warning:*} & --- & Itálico \\
  3 & \texttt{\%FILE\%:\%LINE\%:*: note:*} & Azul & Sublinhado \\
  \bottomrule
\end{tabular}
\end{center}

% === 11. Integração com VSCode ===================================
\section{Integração com VSCode}

\subsection{Extensões recomendadas}

\begin{itemize}[nosep]
  \item \textbf{C/C++} (Microsoft) --- IntelliSense e depuração
  \item \textbf{Arduino} (Microsoft) --- gerenciamento de placas
\end{itemize}

\subsection{IntelliSense para AVR}

Crie o arquivo \texttt{.vscode/c\_cpp\_properties.json} na raiz do seu projeto:

\begin{lstlisting}[language={},caption={.vscode/c\_cpp\_properties.json}]
{
  "configurations": [{
    "name": "AVR",
    "includePath": [
      "${workspaceFolder}/**",
      "C:/ESA/avr_tools/avr8-gnu-toolchain-win32_x86/avr/include",
      "C:/ESA/trampoline/os",
      "C:/ESA/trampoline/machines/avr/arduino/cores/arduino",
      "C:/ESA/trampoline/machines/avr/arduino/libraries/mcp_can/src"
    ],
    "defines": ["F_CPU=16000000", "__AVR_ATmega328P__", "ARDUINO=10813"],
    "compilerPath":
      "C:/ESA/avr_tools/avr8-gnu-toolchain-win32_x86/bin/avr-gcc.exe",
    "cStandard": "c11",
    "cppStandard": "c++14",
    "intelliSenseMode": "gcc-x86"
  }],
  "version": 4
}
\end{lstlisting}

\subsection{Tarefas de build (tasks.json)}

Crie \texttt{.vscode/tasks.json} para automatizar compilação e gravação:

\begin{lstlisting}[language={},caption={.vscode/tasks.json (resumido)}]
{
  "version": "2.0.0",
  "tasks": [
    {
      "label": "Trampoline: goil + make",
      "type": "shell",
      "command": "goil --target=avr/arduino/nano --templates=${input:base}/goil/templates/ ${fileBasename} && python make.py",
      "options": { "cwd": "${fileDirname}" },
      "group": { "kind": "build", "isDefault": true },
      "problemMatcher": ["$gcc"]
    },
    {
      "label": "Trampoline: Upload",
      "type": "shell",
      "command": "arduino-cli upload -b arduino:avr:nano:cpu=atmega328old -p ${input:com} -v -i image.hex -t",
      "options": { "cwd": "${fileDirname}" },
      "problemMatcher": []
    },
    {
      "label": "Arduino: Compile + Upload",
      "type": "shell",
      "command": "arduino-cli compile -b arduino:avr:nano -v -e . && arduino-cli upload -b arduino:avr:nano -p ${input:com} -v -t",
      "group": "build",
      "problemMatcher": ["$gcc"]
    }
  ],
  "inputs": [
    { "id": "com", "description": "COM port", "default": "COM25", "type": "promptString" },
    { "id": "base", "description": "Trampoline base path", "default": "../../../..", "type": "promptString" }
  ]
}
\end{lstlisting}

\begin{dicabox}
Use \textbf{Ctrl+Shift+B} para executar a tarefa de build padrão (``goil + make'') e \textbf{Ctrl+Shift+P $\rightarrow$ Tasks: Run Task} para as demais.
\end{dicabox}

\subsection{Configurações recomendadas}

Adicione ao \texttt{.vscode/settings.json}:

\begin{lstlisting}[language={}]
{
  "files.associations": { "*.oil": "c", "*.ino": "cpp" },
  "C_Cpp.errorSquiggles": "enabled",
  "terminal.integrated.defaultProfile.windows": "Git Bash"
}
\end{lstlisting}

% === 12. Exemplo Prático: LED com Alarme ==========================
\section{Exemplo Prático: LED Controlado por Alarme}

Este exemplo cria uma tarefa no Trampoline que alterna um LED a cada 500\,ms usando um alarme do RTOS. \textbf{Diferente do blink básico}, este exemplo usa duas tarefas com prioridades distintas e comunicação via evento.

\subsection{Estrutura do projeto}

\begin{codigobox}
C:\textbackslash ESA\textbackslash trampoline\textbackslash opt\textbackslash devel\textbackslash led\_alarm\textbackslash\\
\quad led\_alarm.oil\quad\# Descrição OSEK do sistema\\
\quad led\_alarm.cpp\quad\# Código-fonte da aplicação\\
\quad Board.h\quad\quad\quad\quad\# Definições de pinos (copiar da apostila)
\end{codigobox}

\subsection{Arquivo OIL}

\begin{lstlisting}[language=OIL,caption={led\_alarm.oil}]
OIL_VERSION = "2.5" : "led_alarm";

CPU test {
  OS config {
    STATUS = STANDARD;
    BUILD = TRUE {
      TRAMPOLINE_BASE_PATH = "../../../";
      APP_NAME = "image";
      APP_SRC = "led_alarm.cpp";
      CPPCOMPILER = "avr-g++";
      COMPILER = "avr-gcc";
      LINKER = "avr-gcc";
      ASSEMBLER = "avr-gcc";
      COPIER = "avr-objcopy";
      SYSTEM = PYTHON;
      LIBRARY = serial;
    };
    SYSTEM_CALL = TRUE;
  };

  APPMODE stdAppmode {};

  // Alarme que ativa ToggleLed a cada 500ms
  ALARM ledAlarm {
    COUNTER = SystemCounter;
    ACTION = ACTIVATETASK { TASK = ToggleLed; };
    AUTOSTART = TRUE {
      ALARMTIME = 500;
      CYCLETIME = 500;
      APPMODE = stdAppmode;
    };
  };

  // Alarme que ativa PrintStatus a cada 2000ms
  ALARM statusAlarm {
    COUNTER = SystemCounter;
    ACTION = ACTIVATETASK { TASK = PrintStatus; };
    AUTOSTART = TRUE {
      ALARMTIME = 2000;
      CYCLETIME = 2000;
      APPMODE = stdAppmode;
    };
  };

  TASK ToggleLed {
    PRIORITY = 10;
    AUTOSTART = FALSE;
    ACTIVATION = 1;
    SCHEDULE = FULL;
    STACKSIZE = 128;
  };

  TASK PrintStatus {
    PRIORITY = 5;
    AUTOSTART = FALSE;
    ACTIVATION = 1;
    SCHEDULE = FULL;
    STACKSIZE = 256;
  };
};
\end{lstlisting}

\subsection{Código-fonte}

\begin{lstlisting}[language=C++,caption={led\_alarm.cpp}]
#include "tpl_os.h"
#include "Arduino.h"

static bool ledState = false;
static unsigned long toggleCount = 0;

void setup() {
  Serial.begin(115200);
  pinMode(LED_BUILTIN, OUTPUT);
  Serial.println("LED Alarm example started.");
}

// Task de alta prioridade: alterna o LED
TASK(ToggleLed) {
  ledState = !ledState;
  digitalWrite(LED_BUILTIN, ledState ? HIGH : LOW);
  toggleCount++;
}

// Task de baixa prioridade: imprime status na serial
TASK(PrintStatus) {
  Serial.print("LED toggles: ");
  Serial.print(toggleCount);
  Serial.print(" | State: ");
  Serial.println(ledState ? "ON" : "OFF");
}
\end{lstlisting}

\subsection{Compilar e gravar}

Abra um terminal na pasta do projeto e execute:

\begin{codigobox}
cd C:\textbackslash ESA\textbackslash trampoline\textbackslash opt\textbackslash devel\textbackslash led\_alarm\\[4pt]
goil --target=avr/arduino/nano --templates=../../../goil/templates/ led\_alarm.oil\\[4pt]
python make.py\\[4pt]
arduino-cli upload -b arduino:avr:nano:cpu=atmega328old -p COM25 -v -i image.hex -t
\end{codigobox}

Ou use o Remote Flasher:
\begin{enumerate}[nosep]
  \item Copie \texttt{image.hex} para sua máquina
  \item Na aba Flash, selecione o PC, placa e porta
  \item Faça ``Upload + Reset + Flash''
  \item Na aba Serial, abra a porta para ver o status sendo impresso a cada 2 segundos
\end{enumerate}

% === 13. Exemplo: CAN com Trampoline =============================
\section{Exemplo Prático: Comunicação CAN com RTOS}

Este exemplo demonstra duas ECUs trocando mensagens CAN usando tarefas do Trampoline. \textbf{Diferente do exemplo da apostila}, aqui o transmissor envia dados de um potenciômetro e o receptor aciona um LED baseado no valor recebido.

\subsection{ECU Transmissora (can\_pot\_tx)}

\begin{lstlisting}[language=OIL,caption={can\_pot\_tx.oil (trecho relevante)}]
  ALARM sendAlarm {
    COUNTER = SystemCounter;
    ACTION = ACTIVATETASK { TASK = SendPotValue; };
    AUTOSTART = TRUE {
      ALARMTIME = 250;
      CYCLETIME = 250;
      APPMODE = stdAppmode;
    };
  };

  TASK SendPotValue {
    PRIORITY = 20;
    AUTOSTART = FALSE;
    ACTIVATION = 1;
    SCHEDULE = FULL;
    STACKSIZE = 256;
  };
\end{lstlisting}

\begin{lstlisting}[language=C++,caption={can\_pot\_tx.cpp}]
#include "tpl_os.h"
#include "Arduino.h"
#include "Board.h"

#define POT_ID 0x200
MCP_CAN CAN1(CAN1_CS);

void setup() {
  Serial.begin(115200);
  pinMode(ECU1_AIN1, INPUT);
  while (CAN1.begin(MCP_ANY, CAN_500KBPS, MCP_8MHZ) != CAN_OK) {
    delay(200);
  }
  CAN1.setMode(MCP_NORMAL);
  Serial.println("CAN TX: Potentiometer sender ready.");
}

TASK(SendPotValue) {
  int potVal = analogRead(ECU1_AIN1);
  byte data[2] = { highByte(potVal), lowByte(potVal) };
  byte ret = CAN1.sendMsgBuf(POT_ID, CAN_STDID, 2, data);
  if (ret == CAN_OK) {
    Serial.print("TX pot=");
    Serial.println(potVal);
  }
}
\end{lstlisting}

\subsection{ECU Receptora (can\_pot\_rx)}

\begin{lstlisting}[language=C++,caption={can\_pot\_rx.cpp (trecho relevante)}]
TASK(ReceiveAndAct) {
  if (!digitalRead(CAN_INT)) {
    unsigned char dlc = 0;
    unsigned char data[8];
    long unsigned int id;
    CAN1.readMsgBuf(&id, &dlc, data);

    if ((id & CAN_EXTENDED_ID) == POT_ID && dlc >= 2) {
      int potVal = (data[0] << 8) | data[1];
      Serial.print("RX pot=");
      Serial.println(potVal);

      // Aciona LED se pot > 512 (metade da faixa)
      digitalWrite(ECU2_DOUT1, potVal > 512 ? HIGH : LOW);
    }
  }
}
\end{lstlisting}

\begin{dicabox}
Use duas janelas seriais no Remote Flasher (aba Serial, ``+ Add Serial'') para monitorar TX e RX simultaneamente. O autoscroll pode ser desativado para analisar os frames CAN com calma.
\end{dicabox}

% === 14. Configuração dos PCs ====================================
\section{Configuração dos PCs do Laboratório}

\begin{center}
\begin{tabular}{llll}
  \toprule
  \textbf{PC} & \textbf{IP} & \textbf{Método Flash} & \textbf{Placas} \\
  \midrule
  PC 217 & 172.20.36.217 & avrdude  & 4 (com reset) \\
  PC 218 & 172.20.36.218 & avrdude  & 4 (sem reset) \\
  PC 220 & 172.20.36.220 & flash.py & 4 (com reset) \\
  PC 221 & 172.20.36.221 & avrdude  & 4 (sem reset) \\
  \bottomrule
\end{tabular}
\end{center}

Cada placa possui 4 portas ECU (COM), uma porta de reset e um script associado. As portas estão documentadas em \texttt{src/lab\_config.py}.

% === 15. Estrutura do Projeto ====================================
\section{Estrutura do Projeto}

\begin{codigobox}
remote\_flasher/\\
\quad src/\\
\quad\quad main.py\quad\quad\quad\quad\quad\quad\# Ponto de entrada\\
\quad\quad settings.py\quad\quad\quad\quad\# Persistência de configurações\\
\quad\quad workers.py\quad\quad\quad\quad\quad\# Threads (SSH, SCP, SFTP, Serial, Câmera)\\
\quad\quad widgets.py\quad\quad\quad\quad\quad\# LogWidget, StatusIndicator\\
\quad\quad main\_window.py\quad\quad\quad\# MainWindow, CameraPanel\\
\quad\quad lab\_config.py\quad\quad\quad\# Configuração de placas/portas\\
\quad\quad serialterm.py\quad\quad\quad\# Script serial bidirecional (remoto)\\
\quad\quad tabs/\\
\quad\quad\quad vpn\_tab.py\quad\quad\quad\# Aba VPN\\
\quad\quad\quad flash\_tab.py\quad\quad\# Aba Flash\\
\quad\quad\quad serial\_tab.py\quad\quad\# Aba Serial (até 4 conexões)\\
\quad\quad\quad ssh\_tab.py\quad\quad\quad\# Aba SSH + SFTP\\
\quad\quad\quad setup\_tab.py\quad\quad\# Aba Setup + guias IDE\\
\quad assets/icon.ico\\
\quad docs/manual.tex\quad\quad\quad\quad\# Este manual\\
\quad secrets.py\quad\quad\quad\quad\quad\quad\# Credenciais (gitignored)\\
\quad setup\_environment.py\quad\# Script de configuração do ambiente\\
\quad requirements.txt\\
\quad build\_exe.bat
\end{codigobox}

% === 16. Solução de Problemas ====================================
\section{Solução de Problemas}

\subsection{``make.py'' abre diálogo de seleção de programa}
O Windows não associou arquivos \texttt{.py} ao Python. Execute em um \textbf{Prompt de Comando} (CMD) como administrador:
\begin{codigobox}
assoc .py=Python.File\\
ftype Python.File="C:\textbackslash Users\textbackslash \%USERNAME\%\textbackslash AppData\textbackslash Local\textbackslash Programs\textbackslash Python\textbackslash Python311\textbackslash python.exe" "\%1" \%*
\end{codigobox}
Alternativa: sempre use \texttt{python make.py} em vez de \texttt{make.py}.

\subsection{``goil'' não é reconhecido como comando}
Abra um \textbf{novo} terminal após executar o Setup (as variáveis PATH são atualizadas no registro do Windows mas terminais já abertos não as recebem).

\subsection{VPN não conecta}
\begin{itemize}[nosep]
  \item Verifique suas credenciais do CIn
  \item Clique em ``Setup VPN Profile'' para recriar a conexão
  \item Verifique se outra VPN não está ativa
\end{itemize}

\subsection{Erro ao gravar firmware}
\begin{itemize}[nosep]
  \item Verifique se a porta COM está correta
  \item Tente resetar a placa antes de gravar
  \item Verifique se o arquivo \texttt{.hex} é válido
\end{itemize}

\subsection{Terminal serial: linhas quebradas ou caracteres estranhos}
\begin{itemize}[nosep]
  \item Verifique se o baudrate está correto (115200 para Trampoline)
  \item O log agora tem word wrap --- linhas longas (ex: frames CAN) são exibidas corretamente
\end{itemize}

\subsection{Setup: ``Python interpreter not found''}
O executável precisa de um Python instalado no sistema para executar o script de setup. Instale Python 3 e certifique-se de marcar ``Add to PATH'' durante a instalação.

\subsection{Setup: download falha com erro SSL}
Alguns servidores universitários têm certificados com problemas. O script automaticamente tenta novamente sem verificação SSL. Se ainda falhar, baixe o componente manualmente (o log indica a URL).

\end{document}
